\documentclass{article}
\usepackage{graphicx} % Required for inserting images
\usepackage[spanish]{babel}
\usepackage{lmodern}
\renewcommand*{\familydefault}{\sfdefault}
\usepackage[]{geometry}
\usepackage{setspace}
\usepackage{xcolor}
\definecolor{azul}{RGB}{0, 160, 255}
\definecolor{azul2}{RGB}{7, 74, 163}
\usepackage{tikz}
\usepackage{float}
\usepackage{booktabs}
\usepackage{array}
\usepackage{longtable}
\pagenumbering{arabic}
\usepackage[T1]{fontenc}
\usepackage[utf8]{inputenc}
\usepackage{listings}        % bloques de código
\usepackage{hyperref}        % hipervínculos (opcional)
\usepackage{caption}         % personalizar pies de figura/tabla
\usepackage{subcaption} 
% ── Configuración de listings────────────────────────
 \definecolor{codebg}{RGB}{15,23,42}
 \definecolor{codefg}{RGB}{226,232,240}
 \definecolor{codecomment}{RGB}{100,116,139}
 \definecolor{codekeyword}{RGB}{96,165,250}
 \definecolor{codestring}{RGB}{74,222,128}

 \lstdefinestyle{pythonstyle}{
   backgroundcolor=\color{codebg},
   basicstyle=\ttfamily\footnotesize\color{codefg},
   commentstyle=\color{codecomment}\itshape,
   keywordstyle=\color{codekeyword}\bfseries,
   stringstyle=\color{codestring},
   breakatwhitespace=false,
   breaklines=true,
   keepspaces=true,
   numbers=left,
   numberstyle=\tiny\color{codecomment},
   numbersep=8pt,
   showspaces=false,
   showstringspaces=false,
   showtabs=false,
   tabsize=4,
   frame=single,
   rulecolor=\color{codecomment},
   captionpos=b,
   language=Python
}
\lstset{style=pythonstyle}

\usepackage[utf8]{inputenc}
\usepackage[T1]{fontenc}
\usepackage[spanish,es-tabla]{babel}
\usepackage{lmodern}
\usepackage{microtype}
\usepackage[margin=2.5cm]{geometry}
\usepackage{setspace}
\usepackage{graphicx}
\usepackage{float}
\usepackage{tabularx}
\usepackage{array}
\usepackage{xcolor}
\usepackage{listings}
\usepackage{enumitem}
\usepackage[hidelinks]{hyperref}

\setstretch{1.25}
\setlength{\parindent}{0pt}
\setlength{\parskip}{0.55em}
\sloppy
\renewcommand{\arraystretch}{1.25}

\definecolor{codegray}{RGB}{245,245,245}
\definecolor{codeframe}{RGB}{180,180,180}
\definecolor{keywordcolor}{RGB}{120,40,130}
\definecolor{commentcolor}{RGB}{45,110,60}
\definecolor{stringcolor}{RGB}{170,70,30}

\lstdefinestyle{codigo}{
    backgroundcolor=\color{codegray},
    frame=single,
    rulecolor=\color{codeframe},
    basicstyle=\ttfamily\small,
    keywordstyle=\color{keywordcolor}\bfseries,
    commentstyle=\color{commentcolor},
    stringstyle=\color{stringcolor},
    numbers=left,
    numberstyle=\tiny,
    numbersep=8pt,
    breaklines=true,
    breakatwhitespace=false,
    showstringspaces=false,
    tabsize=4,
    columns=fullflexible,
    keepspaces=true,
    captionpos=b
}
\newcommand{\evidencia}[3]{%
  \begin{figure}[H]
    \centering
    \includegraphics[width=0.8\textwidth]{#1}
    \caption{#2}
    \label{fig:#1}
    \nopagebreak\small #3 % Para mostrar el tercer argumento como descripción debajo
  \end{figure}
}


\begin{document}
\setlength{\parindent}{0pt}
\begin{titlepage}
    \thispagestyle{empty}
    \newgeometry{left=2 cm, right=2 cm, top=2 cm, bottom=2 cm}
    \begin{tikzpicture}[overlay, fill opacity= 0.7]
        \rotatebox{-45} {\fill[azul2] (12, -23) rectangle(21, 20);}
        \fill[azul] (12, -27) rectangle(15, 4);
    \end{tikzpicture}
{\huge \textsc{\textbf{Automatización de pruebas y análisis
de calidad de Cal.com}}}\\ [6pt]
    {\Large \textsc{Curso y Grupo: Estándares y Métricas para el Desarrollo de Software, TII05D}}\\[2pt]
    {\Large \textsc{Integrantes:\\Diego Calva \\ Gabriela González\\ Diana Macias\\ Rebeca De La Cruz\\ Yael Cervantes \\https://github.com/UP240080/EyMDSW-Calcom.git}} \\ [2pt]
    \begin{minipage}{3 cm}
    \hspace{0.3 cm} \rotatebox{90}{\Large \textsc{{\color{white}{Ing en}} Tecnologías de la Información e Innovación Digital}}
        
    \end{minipage} \hfill
    \begin{minipage}{4 cm}
        \begin{spacing}{5}
        {\color{white} \fontsize{60}{70} \selectfont \bfseries 2\\ 0\\ 2\\ 6} 
        \end{spacing}
    \end{minipage}
    \vfill
    \begin{flushleft}
      \color{white} { {\large \textsc{Universidad Politécnica\\ de Aguascalientes\\ 21/06/2026}} \\ [10pt]

        }
    \end{flushleft}
\end{titlepage}

\begin{titlepage}
  \Large{\tableofcontents}
 \begin{tikzpicture}[overlay, fill opacity= 0.7]
        \rotatebox{60} {\fill[azul2] (12, -23) rectangle(21, 20);}
        \rotatebox{-85} {\fill[azul2] (12, -23) rectangle(21, 20);}
    \end{tikzpicture}
  
\end{titlepage}

\newpage
\section{Introducción}
El presente trabajo documenta la investigación comparativa y la
implementación práctica de herramientas de automatización de pruebas
y de análisis estático de código, aplicadas al proyecto integrador
del equipo: \textbf{Cal.com}, una plataforma open-source para la
programación de citas y gestión de disponibilidad, construida sobre
Next.js, TypeScript y Prisma, con PostgreSQL como motor de base de
datos. El repositorio oficial sobre el cual se trabaja se encuentra
disponible en \url{https://github.com/UP240080/EyMDSW-Calcom}.

A lo largo del cuatrimestre el equipo ha analizado tres
funcionalidades críticas de la plataforma --gestión de
disponibilidad, flujo de reserva de citas y panel administrativo--
bajo los atributos de calidad de usabilidad, compatibilidad y
accesibilidad. La presente tarea complementa ese análisis al
incorporar dos frentes adicionales: por un lado, una investigación
sobre el estado del arte de los frameworks de pruebas end-to-end
(E2E) y de las herramientas de análisis estático de código; por otro
lado, la aplicación práctica de Selenium con pytest y de SonarQube
directamente sobre el código del proyecto.
Se investigarán al menos tres frameworks de automatización E2E
(Selenium, Cypress y Playwright) y al menos tres herramientas de
análisis estático (SonarQube, ESLint y Semgrep), comparando sus
características técnicas para justificar su pertinencia frente al
stack tecnológico de Cal.com. Posteriormente, el equipo implementará
una suite de pruebas E2E con Page Object Model sobre un módulo real
del flujo de reserva, y ejecutará un análisis estático con
SonarQube para identificar issues de calidad, los cuales serán
priorizados en un plan de remediación de deuda técnica. Finalmente,
se documentará la integración de ambas prácticas en un pipeline de
integración continua.


% ============================================================
% DESARROLLO
% ============================================================

\section{Desarrollo}

Esta sección integra los seis componentes solicitados por la tarea:
investigación de frameworks E2E, investigación de herramientas de
análisis estático, implementación de pruebas con Selenium,
análisis estático con SonarQube, integración en CI y un plan de
remediación de deuda técnica. Las dos primeras subsecciones
(investigación) sientan las bases conceptuales que justifican las
decisiones tomadas en las subsecciones prácticas que les siguen.

% ============================================================
% 4.1
% ============================================================

\subsection{Investigación: frameworks de automatización E2E}

Las pruebas end-to-end (E2E) validan el comportamiento de una
aplicación desde la perspectiva del usuario final, simulando
interacciones reales sobre la interfaz en un navegador. Jorgensen
(2022) ubica este tipo de pruebas en el nivel más alto de la
pirámide de pruebas, ya que ejercitan el sistema integrado de punta
a punta, a diferencia de las pruebas unitarias que aíslan
componentes individuales. Para el caso de Cal.com, cuyo flujo de
reserva de citas involucra múltiples pasos secuenciales
(selección de tipo de evento, elección de fecha y hora, captura de
datos y confirmación), este tipo de prueba resulta indispensable
para validar la experiencia completa del usuario.

Se compararon tres de los frameworks más utilizados actualmente en
la industria: \textbf{Selenium}, \textbf{Cypress} y
\textbf{Playwright}.
Selenium es el framework más longevo de los tres y opera mediante el
protocolo WebDriver, lo que le permite comunicarse con prácticamente
cualquier navegador a través de su respectivo driver.
Su madurez se traduce en una comunidad extensa y en soporte para
múltiples lenguajes de programación, pero también en una curva de
aprendizaje más pronunciada, ya que el manejo de esperas y la
gestión de drivers recae casi por completo en quien escribe la
prueba.
Cypress, en cambio, se ejecuta dentro del mismo bucle de eventos del
navegador, lo que le da acceso directo al DOM y a las peticiones de
red de la aplicación bajo prueba. Esta arquitectura simplifica
considerablemente el manejo de esperas, ya que Cypress reintenta
automáticamente las aserciones hasta que se cumplen o se agota un
tiempo límite. Su principal limitación es que está
acotado al ecosistema JavaScript/TypeScript y, históricamente, a
navegadores basados en Chromium, Firefox y Edge.
Playwright, desarrollado por Microsoft, combina varias de las
ventajas de los dos anteriores: ofrece soporte multinavegador real
(Chromium, Firefox y WebKit) desde un único motor de automatización,
esperas automáticas integradas (auto-waiting) y soporte para
múltiples lenguajes, incluyendo TypeScript, Python, Java y C\#. Su
arquitectura basada en contextos de navegador aislados le permite
ejecutar pruebas en paralelo con mayor velocidad que las
arquitecturas tradicionales basadas en WebDriver (Khorikov, 2020,
respecto a la importancia de la velocidad de ejecución para
mantener suites de prueba sostenibles).

\begin{table}[h]
\centering
\resizebox{\textwidth}{!}{%
\begin{tabular}{lccccccc}
\toprule
\textbf{Criterio} & \textbf{Selenium} & \textbf{Cypress} & \textbf{Playwright} \\
\midrule
Lenguajes soportados & Java, Python, C\#, JS/TS, Ruby, Kotlin & JavaScript / TypeScript & JS/TS, Python, Java, C\# \\
Curva de aprendizaje & Alta (gestión manual de drivers y esperas) & Media (API declarativa, propia de JS) & Media-baja (auto-wait, API moderna) \\
Velocidad de ejecución & Media (overhead del protocolo WebDriver) & Alta (ejecución dentro del navegador) & Alta (contextos aislados, ejecución paralela nativa) \\
Manejo de esperas & Manual (\texttt{WebDriverWait}) & Automático (retry-ability integrado) & Automático (auto-waiting integrado) \\
Soporte de navegadores & Chrome, Firefox, Edge, Safari (vía drivers) & Chrome, Edge, Firefox (Safari experimental) & Chromium, Firefox y WebKit nativos \\
Comunidad & Muy amplia, desde 2004 & Amplia, en crecimiento desde 2017 & En crecimiento acelerado, respaldo de Microsoft \\
Licencia & Apache 2.0 (open source) & MIT (núcleo open source; Cypress Cloud es de pago) & Apache 2.0 (open source) \\
\bottomrule
\end{tabular}%
}
\caption{Comparativa de frameworks de automatización E2E. Fuente: elaboración propia con base en García (2022), Aniche (2022) y documentación oficial (selenium.dev, docs.cypress.io, playwright.dev).}
\label{tab:e2e}
\end{table}

\textbf{Recomendación.} Dado que Cal.com es una aplicación construida
íntegramente en TypeScript sobre Next.js, Playwright sería la
opción más adecuada para un entorno de producción: ofrece soporte
nativo del lenguaje del proyecto, cobertura multinavegador real y
una velocidad de ejecución superior, factores relevantes para una
aplicación con flujos críticos como la reserva de citas. Sin
embargo, y tal como lo permite el lineamiento de la tarea, la
implementación práctica de este trabajo (sección 4.3) se realiza con
\textbf{Selenium + pytest}, por ser la herramienta trabajada en las
sesiones de clase y por contar con un ecosistema robusto de
documentación y patrones (como Page Object Model) ampliamente
documentados por García (2022). Esta decisión no contradice la
recomendación anterior, sino que distingue entre la herramienta
técnicamente más idónea para el stack del proyecto y la herramienta
pedagógicamente adecuada para los objetivos de esta unidad.


% ============================================================
% 4.2
% ============================================================

\subsection{Investigación: análisis estático de código}

El análisis estático de código examina el código fuente sin
ejecutarlo, identificando patrones que indican posibles defectos,
vulnerabilidades de seguridad, código duplicado o violaciones a
buenas prácticas de diseño (Khorikov, 2020). Esto lo distingue del
análisis dinámico --como las pruebas E2E descritas en la sección
anterior-- que observa el comportamiento real del sistema durante su
ejecución y, por lo tanto, solo puede detectar defectos en las rutas
de código que efectivamente se ejercitan durante la prueba. El
análisis estático, en cambio, recorre el cien por ciento del código
fuente, lo que le permite detectar problemas de forma temprana, antes
de que el sistema esté siquiera en ejecución, aunque con el costo de
generar en ocasiones falsos positivos.

Se compararon tres herramientas de análisis estático relevantes para
el stack de Cal.com: \textbf{SonarQube}, \textbf{ESLint} y
\textbf{Semgrep}.

SonarQube es una plataforma integral que evalúa cinco dimensiones de
calidad (confiabilidad, seguridad, mantenibilidad, cobertura y
duplicación) y soporta de forma nativa más de 30 lenguajes,
incluyendo TypeScript y JavaScript. Su fortaleza principal es ofrecer
una visión consolidada del estado de calidad de todo un repositorio,
junto con un mecanismo de \emph{Quality Gate} que puede bloquear
integraciones que no cumplan los umbrales definidos.
ESLint, por su parte, es un linter especializado únicamente en
JavaScript y TypeScript, enfocado en detectar errores de sintaxis,
problemas de estilo y antipatrones específicos del lenguaje en
tiempo de escritura, generalmente integrado directamente en el editor
o como parte de los hooks de pre-commit. A diferencia de SonarQube,
no ofrece una visión histórica ni un tablero centralizado de
calidad, pero su retroalimentación es prácticamente inmediata para
quien programa.
Semgrep se distingue por su enfoque basado en reglas semánticas
ligeras, definidas en patrones legibles similares al propio código,
lo que facilita escribir reglas personalizadas para detectar
vulnerabilidades de seguridad específicas (por ejemplo, uso indebido
de consultas SQL sin parametrizar) de forma más ágil que con motores
de análisis más pesados (Aniche, 2022, respecto a la importancia de
detectar errores comunes de forma temprana y automatizada).

\begin{table}[h]
\centering
\resizebox{\textwidth}{!}{%
\begin{tabular}{lccc}
\toprule
\textbf{Criterio} & \textbf{SonarQube} & \textbf{ESLint} & \textbf{Semgrep} \\
\midrule
Tipo de defectos detectados & Bugs, vulnerabilidades, code smells, duplicación, hotspots de seguridad & Errores de sintaxis, antipatrones de JS/TS, estilo de código & Vulnerabilidades de seguridad y patrones personalizados \\
Lenguajes soportados & +30 lenguajes (incluye TS/JS) & JavaScript / TypeScript (vía plugins) & +20 lenguajes, reglas personalizables \\
Integración en flujo de trabajo & Servidor centralizado / CI, Quality Gate & Editor / pre-commit / CI & CLI / CI, reglas como código \\
Visión histórica del proyecto & Sí (tablero con tendencias) & No & Limitada \\
Licencia & Community Edition gratuita; ediciones avanzadas de pago & MIT (open source) & Open source (núcleo); Semgrep Cloud de pago \\
\bottomrule
\end{tabular}%
}
\caption{Comparativa de herramientas de análisis estático de código. Fuente: elaboración propia con base en Khorikov (2020), Aniche (2022) y documentación oficial (docs.sonarsource.com, eslint.org, semgrep.dev).}
\label{tab:static}
\end{table}

\textbf{Justificación para el stack de Cal.com.} SonarQube resulta
adecuado para este proyecto porque analiza de forma nativa
TypeScript y JavaScript --los lenguajes en los que está escrito
prácticamente todo el código de Cal.com--, y porque su mecanismo de
\emph{Quality Gate} permite establecer un criterio objetivo y
automatizable sobre el cual el equipo puede decidir si un cambio es
apto para integrarse al proyecto, algo especialmente valioso dado el
volumen y la complejidad del repositorio descritos como riesgo en la
elección del caso de estudio. Además, SonarQube puede consumir
directamente el archivo \texttt{coverage.xml} generado por
\texttt{pytest-cov}, lo que permite vincular en un mismo tablero la
cobertura producida por la suite de pruebas E2E (sección 4.3) con el
resto de las métricas estáticas del código (sección 4.4). Por estas
razones, se selecciona SonarQube como herramienta principal de
análisis estático para la práctica, complementando su uso, de ser
necesario, con las reglas de ESLint ya presentes en la configuración
original del proyecto Cal.com.

\subsection{Implementación de pruebas E2E con Selenium}
\label{subsec:selenium}

Esta sección documenta la implementación práctica de una suite de pruebas
\textit{end-to-end} (E2E) sobre el proyecto integrador \textbf{Cal.com},
aplicando el patrón \textit{Page Object Model} (POM) con Selenium~4 y pytest.
La decisión de emplear Selenium obedece a los lineamientos de la unidad,
complementando la recomendación técnica de Playwright presentada en la
sección~4.1: mientras Playwright sería la herramienta óptima para el stack
TypeScript de Cal.com, Selenium con pytest representa la herramienta
pedagógicamente establecida en clase, con amplia documentación en el patrón
POM 

% ------------------------------------------------------------
\subsubsection{Estructura del proyecto de pruebas}
% ------------------------------------------------------------

El proyecto de pruebas se organiza siguiendo la separación de
responsabilidades propuesta por \textcite{Garcia2022}: los
\textit{Page Objects} residen en \texttt{pages/}, las pruebas en
\texttt{tests/}, la configuración global de fixtures en
\texttt{conftest.py} y los reportes generados en \texttt{reports/}.

\begin{figure}[htbp]
  \centering
  \includegraphics[width=0.15\linewidth]{img/Captura de pantalla 2026-06-21 183840.png}
  \caption{Estructura de directorios del proyecto de pruebas E2E de Cal.com.}
  \label{fig:estructura-proyecto}
\end{figure}

La Figura~\ref{fig:estructura-proyecto} muestra la organización del
proyecto. El Cuadro~\ref{tab:archivos-suite} describe el rol de cada
archivo en la suite.

\begin{table}[htbp]
  \centering
  \caption{Archivos del proyecto de pruebas y su responsabilidad.}
  \label{tab:archivos-suite}
  \begin{tabular}{lp{9cm}}
    \toprule
    \textbf{Archivo} & \textbf{Responsabilidad} \\
    \midrule
    \texttt{conftest.py}
      & Fixture global de WebDriver (\texttt{scope=\allowbreak{}function}).
        Crea y destruye una instancia de Chrome por cada prueba, garantizando
        aislamiento total de sesión. Admite el argumento \texttt{-{}-headless}
        para ejecución en CI sin interfaz gráfica. \\
    \texttt{pytest.ini}
      & Configuración global de pytest: rutas de pruebas, marcadores
        personalizados, opciones predeterminadas (\texttt{-v},
        \texttt{-{}-tb=short}) y formato de log. \\
    \texttt{requirements.txt}
      & Dependencias fijadas: \texttt{selenium==4.21.0},
        \texttt{pytest==8.2.2}, \texttt{pytest-html==4.1.1},
        \texttt{pytest-cov==5.0.0}. \\
    \texttt{pages/login\_page.py}
      & Page Object del módulo de autenticación (\texttt{/auth/login}). \\
    \texttt{pages/booking\_page.py}
      & Page Object del flujo de reserva de citas. \\
    \texttt{tests/test\_login.py}
      & Suite de 5 pruebas unitarias + 1 prueba data-driven (7 casos). \\
    \texttt{tests/test\_booking.py}
      & Suite de pruebas del flujo de booking con manejo de skip graceful. \\
    \bottomrule
  \end{tabular}
\end{table}

% ------------------------------------------------------------
\subsubsection{Page Objects implementados}
% ------------------------------------------------------------

Se diseñaron dos \textit{Page Objects} que encapsulan los localizadores
y métodos de negocio de los módulos más críticos de Cal.com
\parencite{Aniche2022}.

\paragraph{LoginPage (\texttt{pages/login\_page.py}).}
Centraliza la interacción con la pantalla de inicio de sesión de Cal.com.
Expone métodos de negocio de alto nivel (\texttt{hacer\_login},
\texttt{login\_exitoso}, \texttt{obtener\_mensaje\_error}) en lugar de
localizadores crudos, respetando el principio de encapsulación del patrón POM.
Todas las esperas se implementan con \texttt{WebDriverWait} y
\texttt{expected\_conditions}; el uso de \texttt{time.sleep} está
explícitamente prohibido en la suite \parencite{Garcia2022}.

El Listado~\ref{lst:login-page} muestra los localizadores y el método
compuesto \texttt{hacer\_login}:

\begin{lstlisting}[
  style=pythonstyle,
  caption={Fragmento de \texttt{pages/login\_page.py}: localizadores y método \texttt{hacer\_login}.},
  label={lst:login-page}
]
# Localizadores centralizados (By, valor)
_CAMPO_EMAIL    = (By.ID, "email")
_CAMPO_PASSWORD = (By.ID, "password")
_BOTON_INGRESAR = (By.CSS_SELECTOR, "button[type='submit']")
_MENSAJE_ERROR  = (By.CSS_SELECTOR,
    "[data-testid='login-error'], .text-red-700, .text-error")
_HEADER_DASHBOARD = (By.CSS_SELECTOR,
    "[data-testid='dashboard'], h3.text-emphasis, nav[aria-label]")

def hacer_login(self, email: str, password: str) -> "LoginPage":
    """Metodo compuesto: llena el formulario y envia en un solo paso."""
    self.ingresar_email(email)
    self.ingresar_password(password)
    self.click_ingresar()
    return self

def _esperar_elemento_visible(self, localizador: tuple):
    """Espera hasta que el elemento sea visible (WebDriverWait)."""
    return self._wait.until(
        EC.visibility_of_element_located(localizador),
        message=f"Elemento no visible: {localizador}",
    )
\end{lstlisting}

\paragraph{BookingPage (\texttt{pages/booking\_page.py}).}
Cubre el flujo completo de reserva de citas: selección de tipo de evento,
elección de fecha en el calendario, selección de horario y llenado del
formulario de confirmación. El método \texttt{abrir\_evento\_directo}
permite saltar al paso del calendario sin pasar por el perfil del usuario,
lo que acorta el tiempo de ejecución de las pruebas relevantes.

El Listado~\ref{lst:booking-page} muestra los localizadores del calendario
y el método de selección de fecha:

\begin{lstlisting}[
  style=pythonstyle,
  caption={Fragmento de \texttt{pages/booking\_page.py}: localizadores del calendario y selección de fecha.},
  label={lst:booking-page}
]
_DIAS_DISPONIBLES  = (By.CSS_SELECTOR,
    "button[data-testid='day']:not([disabled])")
_HORAS_DISPONIBLES = (By.CSS_SELECTOR,
    "button[data-testid='time']")

def seleccionar_primera_fecha_disponible(self) -> "BookingPage":
    """Selecciona el primer dia habilitado del calendario."""
    dias = self.obtener_dias_disponibles()
    if not dias:
        raise AssertionError(
            "No se encontraron fechas disponibles en el calendario.")
    dias[0].click()
    self._esperar_elemento_visible(self._HORAS_DISPONIBLES)
    return self
\end{lstlisting}

% ------------------------------------------------------------
\subsubsection{Suite de pruebas implementadas}
% ------------------------------------------------------------

Se implementaron \textbf{cinco pruebas base} en \texttt{tests/test\_login.py}
que cubren las categorías solicitadas por la rúbrica, más una prueba
data-driven con \texttt{@pytest.mark.parametrize}. El Cuadro~\ref{tab:casos-prueba}
sintetiza los casos.

\begin{table}[htbp]
  \centering
  \caption{Casos de prueba implementados en la suite E2E de Cal.com.}
  \label{tab:casos-prueba}
  \begin{tabular}{llp{5.5cm}l}
    \toprule
    \textbf{ID} & \textbf{Tipo} & \textbf{Descripción} & \textbf{Archivo} \\
    \midrule
    TC-LOGIN-01 & Válido    & Login con credenciales correctas; verifica redirección al dashboard.
                           & \texttt{test\_login.py} \\
    TC-LOGIN-02 & Error     & Login con contraseña incorrecta; verifica mensaje de error.
                           & \texttt{test\_login.py} \\
    TC-LOGIN-03 & Frontera  & Ambos campos vacíos; el sistema no debe avanzar al dashboard.
                           & \texttt{test\_login.py} \\
    TC-LOGIN-04 & Frontera  & Email sin símbolo \texttt{@}; valor justo fuera del dominio válido.
                           & \texttt{test\_login.py} \\
    TC-LOGIN-05 & Error     & Email formalmente válido pero sin cuenta registrada.
                           & \texttt{test\_login.py} \\
    TC-LOGIN-06 & Data-driven & 7 combinaciones de credenciales inválidas con
                               \texttt{@pytest.mark.parametrize}.
                           & \texttt{test\_login.py} \\
    \bottomrule
  \end{tabular}
\end{table}

El diseño de los casos de frontera sigue la técnica de \textit{partición
de equivalencia} y \textit{análisis de valores límite} descrita por
\textcite{Jorgensen2022}: el email sin \texttt{@} representa un valor
justo fuera del dominio de entradas válidas, mientras que la cadena vacía
(\texttt{""}) constituye el límite inferior de longitud del campo.

\paragraph{Prueba data-driven (TC-LOGIN-06).}
El Listado~\ref{lst:data-driven} muestra la implementación de la prueba
parametrizada que cubre siete variaciones de credenciales inválidas en una
sola función:

\begin{lstlisting}[
  style=pythonstyle,
  caption={Prueba data-driven \texttt{TC-LOGIN-06} con \texttt{@pytest.mark.parametrize}.},
  label={lst:data-driven}
]
CREDENCIALES_INVALIDAS = [
    ("",                        "",          "ambos_campos_vacios"),
    ("usuario_sin_arroba.com",  VALID_PWD,   "email_sin_arroba"),
    (VALID_EMAIL,               "a",         "password_de_1_caracter"),
    ("noexiste@calcom-qa.com",  "P@ss2024",  "email_inexistente"),
    ("usuario@",                VALID_PWD,   "email_dominio_faltante"),
    ("@dominio.com",            VALID_PWD,   "email_sin_usuario"),
    (VALID_EMAIL,               "",          "password_vacio"),
]

@pytest.mark.parametrize(
    "email, password, caso",
    CREDENCIALES_INVALIDAS,
    ids=[c[2] for c in CREDENCIALES_INVALIDAS],
)
def test_login_credenciales_invalidas_parametrizado(
        driver, email, password, caso):
    login = LoginPage(driver, base_url=BASE_URL)
    login.abrir()
    login.hacer_login(email, password)
    assert not login.login_exitoso(), (
        f"[{caso}] '{email}' / '{password}' "
        f"no deberia permitir acceso al dashboard."
    )
\end{lstlisting}

% ------------------------------------------------------------
\subsubsection{Uso de esperas explícitas}
% ------------------------------------------------------------

Toda la sincronización entre el driver y la aplicación se implementa
exclusivamente con \texttt{WebDriverWait} y \texttt{expected\_conditions}
de Selenium. El uso de \texttt{time.sleep} está prohibido en la suite
por las razones técnicas señaladas por \textcite{Garcia2022}: las esperas
fijas desperdician tiempo en entornos rápidos y generan fallos intermitentes
en entornos lentos (el problema conocido como \textit{flaky tests}).

Los Page Objects centralizan este comportamiento en dos métodos privados:
\texttt{\_esperar\_elemento\_visible} (usa \texttt{visibility\_of\_element\_located})
y \texttt{\_esperar\_elemento\_clickable} (usa \texttt{element\_to\_be\_clickable}),
con un \textit{timeout} predeterminado de 15 segundos configurable por
instancia. Ninguna prueba llama directamente a \texttt{WebDriverWait};
siempre lo hace a través del Page Object, preservando el encapsulamiento.

% ------------------------------------------------------------
\subsubsection{Ejecución de la suite y resultados}
% ------------------------------------------------------------

La suite se ejecuta con el siguiente comando, que activa simultáneamente
la recolección de cobertura (para posterior consumo por SonarQube) y la
generación del reporte HTML:

\begin{lstlisting}[
  style=pythonstyle,
  language=bash,
  caption={Comando de ejecución de la suite E2E con cobertura y reporte HTML.},
  label={lst:comando-pytest}
]
python -m pytest tests/ ^
    --cov=pages ^
    --cov-report=xml ^
    --html=reports/reporte_e2e.html ^
    --self-contained-html ^
    -v
\end{lstlisting}

\begin{figure}[htbp]
  \centering
  \includegraphics[width=0.65\linewidth]{img/Captura de pantalla 2026-06-21 195221.png}
  \caption{Salida de la terminal tras la ejecución completa de la suite con pytest.}
  \label{fig:terminal-pytest}
\end{figure}

La Figura~\ref{fig:terminal-pytest} muestra la salida de la terminal.
Se ejecutaron 19 casos en total: 12 aprobados (\textit{passed}) y
7 omitido (\textit{skipped}), correspondiente a una prueba del flujo de
booking que requiere la instancia local del proyecto para acceder al
formulario de confirmación.

\begin{figure}[htbp]
  \centering
  \includegraphics[width=0.65\linewidth]{img/Captura de pantalla 2026-06-21 195523.png}
  \caption{Reporte HTML generado por pytest-html con el resumen y detalle de cada caso de prueba.}
  \label{fig:reporte-html}
\end{figure}


El Cuadro~\ref{tab:resultados-suite} consolida los resultados obtenidos.

\begin{table}[htbp]
  \centering
  \caption{Resultados de la ejecución de la suite E2E de Cal.com.}
  \label{tab:resultados-suite}
  \begin{tabular}{lrr}
    \toprule
    \textbf{Módulo / archivo} & \textbf{Casos} & \textbf{Resultado} \\
    \midrule
    \texttt{tests/test\_login.py} — pruebas individuales & 5 & 5 PASSED \\
    \texttt{tests/test\_login.py} — data-driven (TC-LOGIN-06) & 7 & 7 PASSED \\
    \texttt{tests/test\_booking.py} & 7 & 7 SKIPPED \\
    \midrule
    \textbf{Total} & \textbf{19} & \textbf{12 passed, 7 skipped} \\
    \bottomrule
  \end{tabular}
\end{table}

\paragraph{Cobertura de código.}
La ejecución con \texttt{pytest-cov} produjo una cobertura agregada del
\textbf{73\,\%} sobre el directorio \texttt{pages/}:
\texttt{login\_page.py} alcanzó el 88\,\%, mientras que
\texttt{booking\_page.py} llegó al 61\,\% debido a las pruebas de
booking que se ejecutaron en modo \textit{skip} por depender de la
instancia local. El archivo \texttt{coverage.xml} generado se integra
a SonarQube como se describe en la sección~4.4, vinculando en un mismo
tablero la cobertura de pruebas con las métricas de análisis estático.

% ============================================================
%  SECCIÓN 4.4 — Pruebas data-driven (Diego Calva)
% ============================================================
\subsection{Pruebas data-driven y conexión de cobertura con \texttt{pytest-cov}}

\subsubsection{Objetivo de la implementación}

Para reducir la duplicación de código y ejecutar un mismo flujo con distintos datos de entrada, se diseñó una prueba \textit{data-driven} para el módulo de autenticación de Cal.com. La implementación utiliza el decorador \texttt{@pytest.mark.parametrize} y obtiene los datos desde un archivo CSV externo.

La separación entre la lógica de prueba y los datos permite ampliar los escenarios sin crear una función diferente para cada combinación. Este enfoque también mejora la mantenibilidad de la suite y facilita la identificación de fallas específicas.

La cobertura se integró mediante \texttt{pytest-cov}, con salida en terminal, reporte HTML y archivo XML compatible con SonarQube.

\subsubsection{Diseño de los casos de prueba}

La prueba valida el inicio de sesión de Cal.com con datos válidos, inválidos y de frontera. Los escenarios definidos se muestran en la Tabla~\ref{tab:data-driven-login}.

\begin{table}[H]
\centering
\small
\caption{Casos data-driven para el inicio de sesión de Cal.com}
\label{tab:data-driven-login}
\begin{tabularx}{\textwidth}{|>{\centering\arraybackslash}p{1.4cm}|X|X|}
\hline
\textbf{ID} & \textbf{Datos de entrada} & \textbf{Resultado esperado} \\
\hline
DD-01 & Correo y contraseña válidos & El usuario inicia sesión y abandona la ruta \texttt{/auth/login}. \\
\hline
DD-02 & Correo sin formato válido & El navegador muestra una validación del campo de correo. \\
\hline
DD-03 & Correo vacío & El formulario marca el campo de correo como obligatorio. \\
\hline
DD-04 & Contraseña vacía & El formulario marca el campo de contraseña como obligatorio. \\
\hline
DD-05 & Contraseña incorrecta & La aplicación muestra un error de autenticación y no inicia sesión. \\
\hline
\end{tabularx}
\end{table}

Las credenciales válidas se obtienen mediante las variables de entorno \texttt{CALCOM\_TEST\_EMAIL} y \texttt{CALCOM\_TEST\_PASSWORD}. Con ello, el repositorio puede conservar los casos de prueba sin exponer información sensible.

\subsubsection{Archivo de datos CSV}

Los datos de entrada se almacenaron en el archivo \texttt{tests/data/login\_cases.csv}. Cada fila representa una ejecución independiente de la misma prueba.

\begin{lstlisting}[caption={Archivo CSV con los casos de prueba data-driven}]
id,email,password,expected,description
DD-01,ENV:CALCOM_TEST_EMAIL,ENV:CALCOM_TEST_PASSWORD,success,Credenciales validas
DD-02,correo-sin-arroba,Password123!,validation,Formato de correo invalido
DD-03,,Password123!,validation,Correo vacio
DD-04,usuario@example.com,,validation,Contrasena vacia
DD-05,usuario@example.com,ContrasenaIncorrecta123!,auth_error,Contrasena incorrecta
\end{lstlisting}

El prefijo \texttt{ENV:} identifica los valores que se recuperan desde variables de entorno. Esta estructura mantiene separados los datos sensibles y los datos versionados.

\subsubsection{Implementación de la prueba parametrizada}

La función \texttt{load\_login\_cases()} abre el archivo CSV, transforma cada fila en un diccionario y sustituye los valores con prefijo \texttt{ENV:} por el contenido de las variables de entorno correspondientes.

\begin{lstlisting}[language=Python,caption={Prueba data-driven con \texttt{@pytest.mark.parametrize}}]
from __future__ import annotations

import csv
import os
from pathlib import Path
import pytest
from pages.login_page import LoginPage
DATA_FILE = Path(__file__).parent / "data" / "login_cases.csv"

def _resolve_value(value: str) -> str:
    if value.startswith("ENV:"):
        return os.getenv(value.removeprefix("ENV:"), "")
    return value
def load_login_cases() -> list[dict[str, str]]:
    with DATA_FILE.open(encoding="utf-8", newline="") as file:
        cases = list(csv.DictReader(file))
    for case in cases:
        case["email"] = _resolve_value(case["email"])
        case["password"] = _resolve_value(case["password"])
    return cases

@pytest.mark.e2e
@pytest.mark.parametrize(
    "case",
    load_login_cases(),
    ids=lambda case: case["id"],
)
def test_login_data_driven(
    driver,
    base_url: str,
    case: dict[str, str],
) -> None:
    if case["id"] == "DD-01" and (
        not case["email"] or not case["password"]
    ):
        pytest.skip(
            "Credenciales de prueba no disponibles "
            "en las variables de entorno."
        )
    login_page = LoginPage(driver, base_url)
    login_page.open()
    login_page.submit_credentials(
        case["email"],
        case["password"],
    )
    result = login_page.wait_for_result()
    assert result.status == case["expected"], (
        f"Caso {case['id']} ({case['description']}): "
        f"se esperaba '{case['expected']}', "
        f"pero se obtuvo '{result.status}'. "
        f"Detalle: {result.message}"
    )
\end{lstlisting}

Aunque existe una sola función, pytest registra cada fila del archivo CSV como una ejecución separada. Los casos \texttt{DD-01}, \texttt{DD-02}, \texttt{DD-03}, \texttt{DD-04} y \texttt{DD-05} aparecen individualmente en la terminal y en el reporte HTML.

El argumento \texttt{ids=lambda case: case["id"]} asigna a cada ejecución el identificador definido en el CSV. Esto mejora la lectura de los resultados y permite localizar con rapidez el conjunto de datos relacionado con una falla.

\evidencia
{img/salida\_pytest.png}
{Evidencia de ejecución de los casos data-driven en pytest}
{Salida de pytest con la ejecución individual de los casos DD-01 a DD-05.}

\subsubsection{Integración con Page Object Model}

La prueba se apoya en el Page Object \texttt{LoginPage}. Este objeto concentra los localizadores y las operaciones relacionadas con el inicio de sesión: abrir la página, ingresar credenciales, enviar el formulario y esperar el resultado.

\begin{lstlisting}[language=Python,caption={Fragmento principal del Page Object \texttt{LoginPage}}]
from dataclasses import dataclass
from selenium.common.exceptions import TimeoutException
from selenium.webdriver.common.by import By
from selenium.webdriver.support import expected_conditions as EC
from selenium.webdriver.support.ui import WebDriverWait

@dataclass(frozen=True)
class LoginResult:
    status: str
    message: str = ""
class LoginPage:
    EMAIL = (
        By.CSS_SELECTOR,
        'input[name="email"], input[type="email"]',
    )
    PASSWORD = (
        By.CSS_SELECTOR,
        'input[name="password"], input[type="password"]',
    )
    SUBMIT = (By.CSS_SELECTOR, 'button[type="submit"]')
    ERROR = (
        By.CSS_SELECTOR,
        '[role="alert"], [data-testid="toast-error"], '
        '[data-testid*="error"]',
    )

    def __init__(self, driver, base_url: str, timeout: int = 12):
        self.driver = driver
        self.base_url = base_url.rstrip("/")
        self.wait = WebDriverWait(driver, timeout)

    def open(self) -> None:
        self.driver.get(f"{self.base_url}/auth/login")
        self.wait.until(
            EC.visibility_of_element_located(self.EMAIL)
        )

    def submit_credentials(self, email: str, password: str) -> None:
        email_input = self.wait.until(
            EC.visibility_of_element_located(self.EMAIL)
        )
        email_input.clear()
        email_input.send_keys(email)
        password_fields = self.driver.find_elements(*self.PASSWORD)
        if not password_fields:
            self.wait.until(
                EC.element_to_be_clickable(self.SUBMIT)
            ).click()
            password_input = self.wait.until(
                EC.visibility_of_element_located(self.PASSWORD)
            )
        else:
            password_input = password_fields[0]
        password_input.clear()
        password_input.send_keys(password)
        self.wait.until(
            EC.element_to_be_clickable(self.SUBMIT)
        ).click()
\end{lstlisting}

La sincronización con la interfaz se realizó mediante esperas explícitas con \texttt{WebDriverWait}. El uso de condiciones explícitas evita depender de pausas fijas y reduce la inestabilidad de las pruebas.

Los localizadores consideran atributos comunes del formulario de autenticación, como \texttt{name}, \texttt{type}, \texttt{role} y \texttt{data-testid}.

\subsection{Configuración de pytest y pytest-cov}

La suite utiliza Selenium, pytest, pytest-cov y pytest-html. Las dependencias se instalaron con el siguiente comando:

\begin{lstlisting}[caption={Instalación de dependencias para la suite}]
pip install selenium pytest pytest-cov pytest-html
\end{lstlisting}

La configuración de pytest se almacenó en el archivo \texttt{pytest.ini}:

\begin{lstlisting}[caption={Configuración del archivo \texttt{pytest.ini}}]
[pytest]
testpaths = tests
python_files = test_*.py
markers =
    e2e: pruebas de extremo a extremo con navegador
addopts =
    -v
    --strict-markers
    --html=reports/reporte_pytest.html
    --self-contained-html
    --cov=pages
    --cov-branch
    --cov-report=term-missing
    --cov-report=xml:reports/coverage.xml
    --cov-report=html:reports/htmlcov
\end{lstlisting}

Esta configuración genera una salida detallada en terminal, un reporte HTML de la ejecución, un reporte HTML de cobertura y un archivo XML compatible con SonarQube.

La ejecución principal queda reducida al comando siguiente:

\begin{lstlisting}[caption={Ejecución mediante la configuración de \texttt{pytest.ini}}]
pytest
\end{lstlisting}

La misma ejecución puede expresarse de forma completa mediante los parámetros configurados:

\begin{lstlisting}[caption={Ejecución de pruebas y generación de reportes}]
pytest -v \
  --cov=pages \
  --cov-branch \
  --cov-report=term-missing \
  --cov-report=xml:reports/coverage.xml \
  --cov-report=html:reports/htmlcov \
  --html=reports/reporte_pytest.html \
  --self-contained-html
\end{lstlisting}

\subsubsection{Configuración del entorno de prueba}

La instancia local y las credenciales de prueba se administran mediante variables de entorno. En PowerShell se utilizó la siguiente configuración:

\begin{lstlisting}[caption={Variables de entorno en PowerShell}]
$env:CALCOM_BASE_URL="http://localhost:3000"
$env:CALCOM_TEST_EMAIL="usuario_prueba@ejemplo.com"
$env:CALCOM_TEST_PASSWORD="ContrasenaDePrueba123!"
pytest
\end{lstlisting}

La configuración equivalente en Linux o macOS es la siguiente:

\begin{lstlisting}[caption={Variables de entorno en Linux o macOS}]
export CALCOM_BASE_URL="http://localhost:3000"
export CALCOM_TEST_EMAIL="usuario_prueba@ejemplo.com"
export CALCOM_TEST_PASSWORD="ContrasenaDePrueba123!"
pytest
\end{lstlisting}

\subsection{Reportes generados}

La ejecución genera la siguiente estructura de archivos:

\begin{lstlisting}[caption={Archivos generados por pytest}]
reports/
|-- coverage.xml
|-- reporte_pytest.html
`-- htmlcov/
    `-- index.html
\end{lstlisting}

El archivo \texttt{reporte\_pytest.html} contiene el resultado de cada caso de prueba. El archivo \texttt{htmlcov/index.html} presenta la cobertura por archivo y por línea. El archivo \texttt{coverage.xml} contiene la información de cobertura utilizada por SonarQube.

\evidencia
{img/reporte_pytest_html.png}
{Evidencia del reporte HTML de la ejecución}
{Reporte HTML generado por pytest-html.}

\evidencia
{img/reporte_cobertura_html.png}
{Evidencia del reporte HTML de cobertura}
{Reporte de cobertura por archivo y por línea generado por pytest-cov.}

\subsection{Conexión de la cobertura con SonarQube}

La cobertura XML se incorporó al análisis mediante la propiedad \texttt{sonar.python.coverage.reportPaths} del archivo \texttt{sonar-project.properties}:

\begin{lstlisting}[caption={Configuración de cobertura Python en SonarQube}]
sonar.python.coverage.reportPaths=selenium_tests/reports/coverage.xml
\end{lstlisting}

Esta ruta vincula el análisis de SonarQube con el reporte producido por pytest-cov dentro de la carpeta de pruebas Selenium.

\texttt{pytest-cov} mide las líneas de código Python ejecutadas durante la suite. En esta implementación, la medición corresponde principalmente a los Page Objects y a las utilidades utilizadas por Selenium.

Cal.com está desarrollado principalmente con Next.js y TypeScript, por lo que la cobertura generada por pytest-cov no representa por sí sola la cobertura completa de la aplicación. La cobertura del código JavaScript y TypeScript se incorpora mediante un reporte LCOV:

\begin{lstlisting}[caption={Configuración de cobertura JavaScript y TypeScript en SonarQube}]
sonar.javascript.lcov.reportPaths=coverage/lcov.info
\end{lstlisting}

De esta manera, SonarQube puede integrar la cobertura XML del código Python de automatización y la cobertura LCOV del código principal de Cal.com sin confundir ambas mediciones.

\evidencia
{img/sonarqube_cobertura.png}
{Evidencia de la cobertura incorporada en SonarQube}
{Cobertura registrada en el tablero de SonarQube.}

\subsubsection{Interpretación de la implementación}

El diseño data-driven permite verificar múltiples combinaciones de entrada sin duplicar funciones. Los nuevos escenarios se agregan en el archivo CSV, mientras que la lógica central de la prueba permanece sin cambios.

La parametrización también mejora la trazabilidad, ya que cada conjunto de datos conserva un identificador propio dentro de la salida de pytest y del reporte HTML.

El reporte \texttt{term-missing} señala las líneas de Python que no fueron ejecutadas. El reporte HTML facilita la revisión visual por archivo, mientras que el archivo \texttt{coverage.xml} permite trasladar la métrica a SonarQube.

La separación entre la cobertura de Python y la cobertura de TypeScript evita interpretar la ejecución de los Page Objects como si representara la totalidad del código de Cal.com.

\subsubsection{Conclusión}

La implementación integró una prueba data-driven mediante \texttt{@pytest.mark.parametrize}, datos externos en formato CSV y variables de entorno para proteger las credenciales. La estructura permite ejecutar casos válidos, inválidos y de frontera desde una sola función de prueba.

La integración con pytest-cov genera resultados en terminal, un reporte HTML y un archivo XML de cobertura. Este último se conecta con SonarQube mediante \texttt{sonar.python.coverage.reportPaths}.

El enfoque mejora la mantenibilidad, la trazabilidad y la integración de la suite con procesos posteriores de análisis de calidad e integración continua.


% ============================================================
%  SECCIÓN 4.5 — Análisis estático con SonarQube
% ============================================================
\subsection{Análisis estático con SonarQube}

El análisis estático de código constituye una práctica central en la ingeniería de calidad
de software, ya que permite identificar defectos, vulnerabilidades y problemas de
mantenibilidad directamente sobre el código fuente, sin necesidad de ejecutar la
aplicación (Khorikov, 2020). En el contexto del proyecto integrador ---Cal.com, una
plataforma construida sobre Next.js, TypeScript y Prisma--- se optó por SonarQube como
herramienta principal de análisis estático, decisión justificada en la sección de
investigación de análisis estático del presente documento.

\subsubsection{Configuración del proyecto}

El análisis se realizó sobre la instancia local del repositorio disponible en
\url{https://github.com/UP240080/EyMDSW-Calcom}. SonarQube se ejecutó en un
contenedor Docker con la edición \textit{Community}:

\begin{lstlisting}[language=bash,caption={Levantamiento del servidor SonarQube con Docker}]
docker run -d --name sonarqube \
  -p 9000:9000 \
  sonarqube:community
\end{lstlisting}

El análisis del código fuente se lanzó mediante el escáner oficial de SonarSource:

\begin{lstlisting}[language=bash,caption={Ejecución del sonar-scanner sobre el repositorio}]
docker run --rm \
  -v "$PWD:/usr/src" \
  --network host \
  sonarsource/sonar-scanner-cli \
  -Dsonar.projectKey=calcom \
  -Dsonar.sources=apps/web/pages,packages \
  -Dsonar.exclusions=**/node_modules/**,**/*.test.ts,**/*.spec.ts \
  -Dsonar.host.url=http://localhost:9000 \
  -Dsonar.login=admin \
  -Dsonar.password=admin
\end{lstlisting}

El archivo \texttt{sonar-project.properties} centraliza los parámetros de configuración
para facilitar la integración con el pipeline de CI:

\begin{lstlisting}[caption={Archivo sonar-project.properties del proyecto}]
sonar.projectKey=calcom-integrador
sonar.projectName=Cal.com - EyMDSW
sonar.projectVersion=1.0
sonar.sources=apps/web/pages,packages
sonar.tests=selenium_tests/tests
sonar.exclusions=**/node_modules/**,**/*.test.ts,**/*.spec.ts,**/.next/**
sonar.language=ts
sonar.python.coverage.reportPaths=selenium_tests/reports/coverage.xml
sonar.javascript.lcov.reportPaths=coverage/lcov.info
sonar.host.url=http://localhost:9000
sonar.login=admin
\end{lstlisting}

\subsubsection{Tablero de calidad del proyecto}

Una vez ejecutado el escáner, el tablero de SonarQube consolidó el estado de calidad
del repositorio bajo las cinco dimensiones definidas por la plataforma.

\evidencia
{img/sonar_dashboard.png}
{Tablero general de SonarQube para Cal.com}
{Tablero general de SonarQube para el proyecto Cal.com.}

\subsubsection{Las cinco dimensiones de calidad}

SonarQube evalúa la calidad del código a través de cinco dimensiones principales.
La Tabla~\ref{tab:sonar_dimensiones} resume los resultados obtenidos para el proyecto Cal.com.

\begin{longtable}{|p{3.2cm}|p{1.8cm}|p{7.5cm}|}
\hline
\textbf{Dimensión} & \textbf{Rating} & \textbf{Descripción del resultado} \\
\hline
\endfirsthead
\hline
\textbf{Dimensión} & \textbf{Rating} & \textbf{Descripción del resultado} \\
\hline
\endhead

Reliability & B & Se detectaron \textit{bugs} clasificados como mayores que podrían provocar comportamiento inesperado en el flujo de reservas. El rating A exige cero bugs. \\
\hline
Security & A & No se reportaron vulnerabilidades críticas en el código analizado. Sin embargo, se identificaron \textit{security hotspots} que requieren revisión manual. \\
\hline
Maintainability & A & La deuda técnica acumulada se estimó en menos del 5\,\% del tiempo de desarrollo total del proyecto, dentro del umbral aceptable. \\
\hline
Coverage & -- & La cobertura reportada corresponde a los Page Objects de la suite Selenium (89.33\,\%, generada por \texttt{pytest-cov}). La cobertura del código TypeScript requiere un reporte LCOV generado por las pruebas nativas del proyecto. \\
\hline
Duplications & -- & Se identificó duplicación de código en componentes de UI del directorio \texttt{apps/web/pages}, principalmente en patrones de manejo de formularios y llamadas a la API. \\
\hline
\caption{Resumen de las cinco dimensiones de calidad en SonarQube para Cal.com.}
\label{tab:sonar_dimensiones}
\end{longtable}

\paragraph{Reliability.}
Los \textit{bugs} detectados se concentraron en lógica de manejo de valores nulos
en componentes de disponibilidad y en llamadas a \texttt{Promise} sin manejo de rechazo.

\paragraph{Security.}
Los \textit{security hotspots} son fragmentos sensibles que no representan
vulnerabilidades confirmadas, pero requieren revisión humana para descartar un riesgo
real (SonarSource, s.f.). En Cal.com, los hotspots detectados involucran el uso de
\texttt{localStorage} y la exposición de información en mensajes de error.

\paragraph{Maintainability.}
La deuda técnica en SonarQube se expresa como el tiempo estimado necesario para
refactorizar todos los \textit{code smells}. En el proyecto, la mayoría correspondían a
funciones con demasiados parámetros, variables no utilizadas y bloques con complejidad
ciclomática elevada.

\paragraph{Coverage.}
La cobertura configurada vincula el reporte \texttt{coverage.xml} generado por
\texttt{pytest-cov} con el código Python de los Page Objects, no con el código
TypeScript de la aplicación principal.

\paragraph{Duplications.}
SonarQube marcó como duplicados bloques de código con alta similitud estructural en
los componentes de formularios del módulo de reservas. La duplicación aumenta el
riesgo de inconsistencias cuando se realizan cambios, ya que los bloques copiados
deben actualizarse por separado (Aniche, 2022).

\subsubsection{Análisis de issues reales encontrados}

Se seleccionaron tres \textit{issues} representativos del análisis, cubriendo distintos
tipos y niveles de severidad.

\paragraph{Issue 1 --- Bug mayor: Promise sin manejo de rechazo.}
\begin{itemize}[leftmargin=*]
  \item \textbf{Tipo:} Bug \quad \textbf{Severidad:} Major
  \item \textbf{Regla:} \texttt{typescript:S4822}
  \item \textbf{Archivo:} \texttt{apps/web/pages/api/availability/index.ts}, línea 47
  \item \textbf{Descripción:} Una llamada asíncrona a la base de datos mediante Prisma
        se realiza sin encadenar un bloque \texttt{.catch()} ni envolverla en
        \texttt{try-catch}. Si la consulta falla, la promesa rechazada no se maneja,
        lo que puede provocar un error no capturado que detenga el servidor Node.js
        o devuelva una respuesta HTTP 500 sin información útil al cliente.
  \item \textbf{Impacto:} Alto. El flujo de reservas depende de consultas de
        disponibilidad; un error no capturado aquí impide completar una cita.
\end{itemize}

\paragraph{Issue 2 --- Code Smell: función con complejidad ciclomática excesiva.}
\begin{itemize}[leftmargin=*]
  \item \textbf{Tipo:} Code Smell \quad \textbf{Severidad:} Critical
  \item \textbf{Regla:} \texttt{typescript:S3776}
  \item \textbf{Archivo:} \texttt{packages/core/EventManager.ts}, línea 112
  \item \textbf{Descripción:} La función \texttt{buildEventData()} presenta una
        complejidad ciclomática de 18, superando el umbral recomendado de 10
        (Jorgensen, 2022). Concatena múltiples condiciones \texttt{if}/\texttt{else}
        para determinar el comportamiento según el tipo de evento, lo que la hace
        difícil de probar y de modificar sin introducir regresiones.
  \item \textbf{Impacto:} Alto sobre la mantenibilidad. Cualquier modificación requiere
        comprender y recorrer toda la función.
\end{itemize}

\paragraph{Issue 3 --- Security Hotspot: uso de \texttt{localStorage} para datos de sesión.}
\begin{itemize}[leftmargin=*]
  \item \textbf{Tipo:} Security Hotspot \quad \textbf{Severidad:} High (pendiente de revisión)
  \item \textbf{Regla:} \texttt{typescript:S5122}
  \item \textbf{Archivo:} \texttt{apps/web/lib/auth/session.ts}, línea 23
  \item \textbf{Descripción:} Se almacena un token de autenticación en
        \texttt{localStorage}, accesible desde cualquier script JavaScript en la
        misma página. Esto es susceptible a ataques XSS si algún componente de
        terceros introduce código malicioso.
  \item \textbf{Impacto:} Potencialmente crítico en entornos de producción. Requiere
        revisión manual para confirmar si el riesgo está mitigado por otros controles.
\end{itemize}

\subsubsection{Conexión de la cobertura de pytest con SonarQube}

La cobertura generada por la suite Selenium se conectó a SonarQube mediante la
propiedad \texttt{sonar.python.coverage.reportPaths} apuntando a
\texttt{selenium\_tests/reports/coverage.xml}, producido por \texttt{pytest-cov}.

De manera complementaria, la cobertura del código TypeScript se incorpora mediante
\texttt{sonar.javascript.lcov.reportPaths}, consumiendo el reporte LCOV generado por
las pruebas nativas del proyecto. Esta separación evita que SonarQube confunda la
cobertura de los scripts Python de automatización con la del código de producción.

\evidencia
{img/sonar_issues.png}
{Lista de issues detectados por SonarQube en Cal.com}
{Issues detectados por SonarQube en el análisis del repositorio Cal.com.}


% ============================================================
%  SECCIÓN 4.6 — Plan de remediación de deuda técnica
% ============================================================
\subsection{Plan de remediación de deuda técnica}

La deuda técnica es la diferencia entre el estado actual del código y el estado ideal
que permitiría mantenerlo, extenderlo y probarlo con el menor esfuerzo posible
(Khorikov, 2020). SonarQube cuantifica esta deuda como el tiempo total estimado para
corregir todos los \textit{code smells} del proyecto; sin embargo, no todos los issues
tienen el mismo valor de remediación: algunos afectan directamente la confiabilidad
del sistema, otros la seguridad, y otros solo impactan la legibilidad del código.

El enfoque \textbf{Clean as You Code} ---propuesto por SonarSource como estrategia de
adopción gradual--- establece que el equipo debe priorizar la calidad del código nuevo
por encima de la remediación masiva del código heredado: cada cambio que se incorpore
al repositorio debe cumplir los umbrales del Quality Gate antes de fusionarse
(SonarSource, s.f.). Bajo este criterio, la priorización de la tabla siguiente
privilegia los issues que afectan los módulos que el equipo ha modificado o que
modificará en las próximas iteraciones del proyecto integrador.

\subsubsection{Criterios de priorización}

Los issues se priorizaron según los siguientes criterios, aplicados en orden:

\begin{enumerate}[leftmargin=*]
  \item \textbf{Tipo de issue:} Los \textit{bugs} tienen precedencia sobre los
        \textit{hotspots} de seguridad, y estos sobre los \textit{code smells},
        porque los bugs representan defectos con alta probabilidad de manifestarse
        en tiempo de ejecución (Jorgensen, 2022).
  \item \textbf{Severidad asignada por SonarQube:} Critical $>$ Major $>$ Minor.
  \item \textbf{Impacto en flujos críticos:} Los módulos de reserva y autenticación
        reciben mayor prioridad por ser los flujos que el equipo ha analizado y
        probado a lo largo de la Unidad II.
  \item \textbf{Esfuerzo estimado vs.\ impacto:} Issues de alto impacto y bajo
        esfuerzo se ubican por encima de issues de menor impacto aunque sean
        fáciles de corregir.
\end{enumerate}

\subsubsection{Tabla de remediación priorizada}

\begin{longtable}{|p{0.5cm}|p{4.2cm}|p{2.1cm}|p{1.5cm}|p{1.5cm}|p{2.1cm}|}
\hline
\textbf{\#} & \textbf{Issue} & \textbf{Tipo / Severidad} & \textbf{Esfuerzo} & \textbf{Prioridad} & \textbf{Responsable} \\
\hline
\endfirsthead
\hline
\textbf{\#} & \textbf{Issue} & \textbf{Tipo / Severidad} & \textbf{Esfuerzo} & \textbf{Prioridad} & \textbf{Responsable} \\
\hline
\endhead

1 & Promise sin manejo de rechazo en API de disponibilidad (\texttt{S4822}) & Bug / Major & 30 min & \textbf{Alta} & Backend \\
\hline
2 & Función \texttt{buildEventData()} con complejidad ciclomática~18 (\texttt{S3776}) & Smell / Critical & 3 h & \textbf{Alta} & Backend \\
\hline
3 & Token de sesión en \texttt{localStorage} (\texttt{S5122}) & Hotspot / High & 2 h & \textbf{Alta} & Frontend \\
\hline
4 & Variables declaradas sin uso en módulo de reservas (\texttt{S1481}) & Smell / Minor & 15 min & \textbf{Media} & Todo el equipo \\
\hline
5 & Bloques duplicados en formularios de reserva & Smell / Major & 2 h & \textbf{Media} & Frontend \\
\hline
6 & Manejo de errores ausente en proveedor de calendario externo & Bug / Minor & 1 h & \textbf{Media} & Backend \\
\hline
7 & Funciones con demasiados parámetros en utilidades de Prisma (\texttt{S107}) & Smell / Minor & 1 h & \textbf{Baja} & Equipo (backlog) \\
\hline
8 & Comentarios \texttt{TODO} sin ticket en código de producción & Smell / Info & 10 min & \textbf{Baja} & Todo el equipo \\
\hline
\caption{Plan de remediación de deuda técnica priorizado por severidad e impacto}
\label{tab:remediation}
\end{longtable}

\subsubsection{Justificación del orden}

\paragraph{Issues 1 y 2 (Alta).}
La promesa sin manejo de rechazo (issue~1) puede provocar la caída del servidor
Node.js en tiempo de ejecución, afectando directamente la disponibilidad de Cal.com.
Su corrección es inmediata: basta con agregar un bloque \texttt{try-catch} alrededor
de la consulta Prisma o encadenar un \texttt{.catch()}.

El issue~2 ---la función con complejidad ciclomática de 18--- no produce un fallo
inmediato, pero representa un riesgo de regresión elevado: cualquier nueva condición
añadida a \texttt{buildEventData()} tiene alta probabilidad de introducir un defecto
en una rama existente no cubierta por pruebas (Aniche, 2022). Reducir la complejidad
extrayendo funciones auxiliares por tipo de evento es una inversión que reducirá el
costo de las futuras iteraciones del proyecto integrador.

\paragraph{Issue 3 (Alta --- Seguridad).}
El almacenamiento de tokens en \texttt{localStorage} es una vulnerabilidad
ampliamente documentada en aplicaciones web. Aunque la explotación requiere un
vector XSS previo, Cal.com integra componentes de calendario de terceros que
representan una superficie de ataque no trivial. La migración hacia cookies
\texttt{HttpOnly} es un cambio de dos horas con impacto de seguridad significativo
y sin efectos secundarios observables para el usuario final.

\paragraph{Issues 4, 5 y 6 (Media).}
Estos issues no afectan el funcionamiento actual del sistema, pero incrementan el
esfuerzo de mantenimiento. La duplicación de código en formularios (issue~5) es
especialmente relevante porque los tests E2E de la sección anterior ya cubren esos
formularios: si el equipo modifica uno de los bloques duplicados pero no el otro,
las pruebas podrían no detectar la inconsistencia.

\paragraph{Issues 7 y 8 (Baja).}
Son mejoras de higiene que se abordan de forma continua. Al abrir un Pull Request
que toque los archivos afectados, el desarrollador responsable corrige el smell
antes de solicitar la revisión. Esta práctica evita que la deuda técnica baja se
acumule sin requerir una sesión de remediación dedicada.

\subsubsection{Conexión con Clean as You Code}

El enfoque \textit{Clean as You Code} propone que el Quality Gate evalúe únicamente
el código nuevo o modificado en cada análisis, de modo que el equipo no se bloquee
por deuda acumulada en código heredado, pero sí garantice que las nuevas
contribuciones elevan el estándar de calidad (SonarSource, s.f.). Para Cal.com se
propone el siguiente Quality Gate mínimo:

\begin{itemize}[leftmargin=*]
  \item Cero bugs nuevos de severidad Major o Critical.
  \item Cero vulnerabilidades de seguridad confirmadas.
  \item Cobertura del código nuevo $\geq$ 80\,\%.
  \item Duplicación de código nuevo $\leq$ 3\,\%.
  \item Todos los Security Hotspots revisados (estado \textit{Reviewed}).
\end{itemize}

Este criterio es coherente con las condiciones de bloqueo del pipeline de CI: un
análisis que no cumpla estos umbrales impide la fusión automática de la rama al
repositorio principal.

\evidencia
{img/sonar_quality_gate.png}
{Quality Gate configurado en SonarQube para el proyecto Cal.com}
{Quality Gate del proyecto Cal.com en SonarQube.}

\subsection{Conclusión integral}
La implementación de herramientas de automatización y análisis estático en nuestra plataforma J\&C representó un punto de inflexión en nuestra forma de concebir la calidad del software. Al transicionar de pruebas manuales a una suite E2E estructurada con Selenium y el patrón Page Object Model, comprendimos el valor de la mantenibilidad en el código de pruebas. Asimismo, la integración de SonarQube nos expuso a la realidad de nuestra deuda técnica; dimensionar los defectos y vulnerabilidades nos demostró que no basta con que el código funcione, sino que debe ser seguro y confiable. Finalmente, consolidar este flujo en un entorno de versionamiento colaborativo con Git fortaleció nuestra dinámica de trabajo, evidenciando que la calidad es un esfuerzo continuo y automatizable.

\subsection{Aportes individuales y control de versiones}
El desarrollo de esta tarea se gestionó íntegramente mediante control de versiones con Git. El historial de \textit{commits} en el repositorio del proyecto evidencia la participación colaborativa de todos los integrantes. A continuación, se detalla la contribución específica de cada miembro:

\begin{itemize}
    \item \textbf{Yael Cervantes:} Lideró la fase de investigación teórica, comparando exhaustivamente tres frameworks de automatización de pruebas E2E y tres herramientas de análisis estático de código. Su análisis crítico permitió fundamentar las decisiones tecnológicas del equipo.
    
    \item \textbf{Gabriela González:} Estuvo a cargo de la implementación principal de la suite de pruebas E2E utilizando Selenium y \texttt{pytest}. Diseñó la arquitectura de pruebas aplicando el patrón Page Object Model (POM) para asegurar la escalabilidad.
    
    \item \textbf{Diego Calva:} Complementó la suite de pruebas mediante el diseño e implementación de pruebas \textit{data-driven}. Además, fue responsable de configurar y conectar la cobertura de código utilizando \texttt{pytest-cov}, asegurando la medición de las pruebas.
    
    \item \textbf{Rebeca De La Cruz:} Dirigió el análisis de calidad del proyecto integrador configurando e interpretando los resultados en SonarQube. A partir de las dimensiones evaluadas, estructuró el plan priorizado de remediación de deuda técnica.
    
    \item \textbf{Diana Macias:} Lideró la consolidación documental del proyecto y la gestión del control de versiones. Se encargó de redactar la conclusión integral, documentar el aporte de cada integrante, administrar el repositorio en Git y asegurar el cumplimiento de la estructura formal del documento.
\end{itemize}

% ============================================================
%  REFERENCIAS
% ============================================================
\section*{Referencias}
\addcontentsline{toc}{section}{Referencias}

Aniche, M. (2022). \textit{Effective software testing: A developer's guide}. Manning.

García, B. (2022). \textit{Hands-on Selenium WebDriver with Java}. O'Reilly Media.

Jorgensen, P. C. (2022). \textit{Software testing: A craftsman's approach} (5.a ed.). CRC Press.

Khorikov, V. (2020). \textit{Unit testing: Principles, practices, and patterns}. Manning.

pytest. (s.f.). \textit{How to parametrize fixtures and test functions}. \url{https://docs.pytest.org}

pytest-cov. (s.f.). \textit{Coverage reporting with pytest-cov}. \url{https://pytest-cov.readthedocs.io}

SonarSource. (s.f.). \textit{Python test coverage}. Documentación oficial de SonarQube. \url{https://docs.sonarsource.com}





\end{document}
